\section{Plumbing geometry and component potentials}\label{sec:plumbing}

The smooth total-space and nodal-fibre assumptions give holomorphic
coordinates $u,v$ at each node in which $\pi=uv$. Choose disjoint
such bidiscs and a fixed small outer radius $r_*>0$. On the two
central branches, $u$ and $v$ are puncture coordinates. We use the
same coordinates to normalize the logarithmic parts of Green functions.

\begin{lemma}\label{lem:plumbing}
There are identifications of the thick part of $X_t$ with the
components punctured at radius $\rho$, for $|t|\ll\rho^2$ and
$0<\rho\ll r_*$, having the following properties.
\begin{enumerate}
\item In the shrinking branch collars they are exactly
$u\mapsto(u,t/u)$ or $v\mapsto(t/v,v)$. On a fixed exterior they
depend smoothly on the real and imaginary parts of $t$ and equal
the identity at $t=0$.
\item In each fixed derivative order $k$, pulled-back metric and
complex-structure coefficients differ from the central ones by at
most $C_k|t|\rho^{-n_k}$ for some fixed integer $n_k$.
\item The middle part of each neck, where $|u|,|v|\le2\rho$, has
area at most $C\rho^2$. Each transition annulus $|u|\asymp\rho$
has area $O(\rho^2)$ and metric uniformly comparable to that of the
central branch.
\end{enumerate}
All constants are uniform in the argument of $t$.
\end{lemma}
\begin{proof}
Remove fixed small node discs from the central components. The map
$\pi$ is a submersion on a compact neighbourhood of the resulting
exterior, including its boundary. Choose smooth horizontal lifts
of the two real parameter directions. On a collar of each boundary,
choose the lifts that keep $u$, or respectively $v$, fixed. A partition
of unity interpolates these lifts only on fixed annuli, where all
coefficients are smooth and bounded. Flow along the straight segment
from zero to $t$. This gives a smooth trivialization of the fixed
exterior which agrees with the coordinate identification on smaller
boundary collars. Extend it inward by $u\mapsto(u,t/u)$ on each
branch. The two branch images are disjoint when $|t|\ll\rho^2$.

On the fixed exterior all coefficient differences are $O_k(|t|)$.
On the shrinking collar,
\[
 |D^k(t/u)|\le C_k|t|\rho^{-k-1}.
\]
The induced metric is the restriction of a smooth ambient tensor
along this map. Differentiating it gives a finite power of
$\rho^{-1}$ in every fixed order. The complex coordinate on the
collar is exactly $u$, so its complex structure needs no correction
there. The smooth branch metric is nondegenerate through the
puncture, and the same estimates hold for inverse coefficients.

The ambient metric is bounded above by a fixed multiple of the flat
metric in the bidisc. In the parameter $u$, its fibre area is bounded by
\[
 C\int_{|t|/(2\rho)\le|u|\le2\rho}
       (1+|t|^2/|u|^4)\,dx\,dy\le C'\rho^2.
\]
On $|u|\asymp\rho$ the second term is uniformly small, and the
induced metric is comparable to the branch metric. The estimates
involve only $|t|$.
\end{proof}

We shall take
\begin{equation}\label{eq:cutoff-radius}
 \rho=L^{-1}.
\end{equation}
Every geometric coefficient error required below is then bounded
by $C e^{-L}L^n$, and is smaller than any negative power of $L$.
Fixed compact subsets are transported by the exterior identifications
from Lemma~\ref{lem:plumbing}.

\begin{lemma}\label{lem:green}
For each node puncture $p\in C_i$, there is a real mean-zero function
$g_{i,p}$, smooth off $p$, such that
\begin{equation}\label{eq:green}
 \Delta_i g_{i,p}=-2\pi\delta_p+\frac{2\pi}{A_i},
 \qquad\int_{C_i}g_{i,p}\,dA_i=0.
\end{equation}
In the assigned puncture coordinate $u$, it has the form
$g_{i,p}=\log|u|+\gamma_{i,p}(u)$, with $\gamma_{i,p}$ smooth.
\end{lemma}
\begin{proof}
For the nonnegative Laplacian,
$\Delta\log|u|=-2\pi\delta_0$. Multiply the logarithm by a cutoff
equal to one near the puncture. Its Laplacian is the stated point
mass plus a smooth function. The right side of \eqref{eq:green}
has integral zero, so the difference is a smooth mean-zero function.
Invert the scalar Laplacian on the mean-zero functions of the fixed
compact smooth component and adjust the additive constant. Elliptic
regularity makes the correction smooth. This also proves uniqueness.
\end{proof}

Along the finite trajectory let
$B_{ij}(\tau)=\log(H_i(\tau)^{-1}H_j(\tau))$, and define
\begin{equation}\label{eq:component-potential}
 \xi_i(\tau,z)=-\sum_{p:\,i\leftrightarrow j}B_{ij}(\tau)g_{i,p}(z),
 \qquad V_i=H_i^{-1}\partial_\tau H_i.
\end{equation}
These matrices are $H_i$-self-adjoint and trace free. The essential
compatibility condition is
\begin{equation}\label{eq:compatibility}
 \Delta_i\xi_i=-V_i\qquad\text{away from the punctures}.
\end{equation}
Indeed \eqref{eq:green} gives the left side as
$-(2\pi/A_i)\sum B_{ij}$, exactly the negative of
\eqref{eq:finite-flow}.

For $\tau\in[0,T]$, the finite trajectory and its derivatives are
bounded. Consequently, away from discs of radius $\rho$,
\begin{equation}\label{eq:potential-bounds}
 \begin{gathered}
 |\xi_i|\le C_T(1+|\log\rho|),\qquad
 |\nabla\xi_i|\le C_T(1+r^{-1}),\\
 |\nabla^2\xi_i|\le C_T(1+r^{-2}),\qquad
 \int|\xi_i|\,dA_i\le C_T,\qquad
 \int|\nabla\xi_i|^2\,dA_i\le C_T(1+|\log\rho|),
 \end{gathered}
\end{equation}
where $r$ is distance in a puncture chart when near a puncture.
The same bounds hold after any fixed number of slow-time derivatives.
They follow by integrating $|\log r|$ and $r^{-2}$ against the
smooth component area. All norms of matrices in this proof are
equivalent on the compact range of the component metrics.
